% GENERATED FILE. Edit canonical lessons, then run npm run generate:books.
% canonical-source-hash: fnv1a64:d97daebeb6674813
% canonical-lessons: MR-C542-niyamavali, MR-C542-niyantran, MR-C542-takid, MR-C542-jumanne, MR-C542-jhugarne

\chapter{Rules and Defiance}
\label{ch:mr-rules-and-defiance}
\hlchaptermodality{542}

\begin{chapteropening}
\textbf{By the end of this chapter:} \emph{I can talk about rulebooks, warnings and defying rules.}

Complete the last lesson of chapter 542: I can talk about rulebooks, warnings and defying rules.
\end{chapteropening}

\section[\texorpdfstring{\mr{नियमावली} (niyamāvalī)}{नियमावली (niyamāvalī)}]{\mr{नियमावली} (niyamāvalī) — a rulebook}

\label{lesson:MR-C542-niyamavali}



\begin{quote}
Before the new one: say the Marathi for alert, then the Marathi for to issue a warning.
\end{quote}

\subsection*{You'll want to know: \mr{नियमावली}}
\textbf{\mr{नियमावली}} — \emph{niyamāvalī} — \textquotedblleft{}a rulebook\textquotedblright{}.

\begin{quote}
\mr{शाळेने} \mr{सगळ्या} \mr{पालकांना} \mr{नियमावली} \mr{दिली}.
\end{quote}

\begin{quote}
\emph{śāḷene sagḷyā pālkānnā niyamāvalī dilī.}
\end{quote}

\begin{quote}
The school gave all the parents a rulebook.
\end{quote}

Feminine.

\begin{cousinweb}
\textbf{\mr{नियमावली}} (\emph{niyamāvalī}) is a tatsama compound of Sanskrit \emph{niyama}, \textquotedblleft{}rule\textquotedblright{}, and \emph{āvali}, \textquotedblleft{}a row\textquotedblright{} (Monier-Williams): a row of rules.
\end{cousinweb}

\begin{grammarlens}[title={What you've built}]
One more word for saying what happens if.
\end{grammarlens}

\subsection*{Your turn}
\begin{itemize}
\raggedright
  \item \emph{Say it:} \emph{niyamāvalī}
  \item \emph{Say it:} \emph{niyamāvalī}, once more
  \item \emph{Say it:} \emph{śāḷene sagḷyā pālkānnā niyamāvalī dilī.} — The school gave all the parents a rulebook.
  \item \emph{Recall:} say the Marathi for alert, then the Marathi for to issue a warning, then say \emph{niyamāvalī} again
\end{itemize}

\begin{tcolorbox}[breakable,colback=teal!4,colframe=teal!35!black,title={Before you move on}]
What does \mr{नियमावली} mean? (\textquotedblleft{}A rulebook\textquotedblright{}.) Say it once more.
\end{tcolorbox}
\section[\texorpdfstring{\mr{नियंत्रण} (niyantraṇ)}{नियंत्रण (niyantraṇ)}]{\mr{नियंत्रण} (niyantraṇ) — control}

\label{lesson:MR-C542-niyantran}



\begin{quote}
Before the new one: say the Marathi for to issue a warning, then the Marathi for a rulebook.
\end{quote}

\subsection*{You'll want to know: \mr{नियंत्रण}}
\textbf{\mr{नियंत्रण}} — \emph{niyantraṇ} — \textquotedblleft{}control\textquotedblright{}.

\begin{quote}
\mr{आग} \mr{लवकरच} \mr{नियंत्रणात} \mr{आली}.
\end{quote}

\begin{quote}
\emph{āg lavkarac niyantraṇāt ālī.}
\end{quote}

\begin{quote}
The fire was soon brought under control.
\end{quote}

Neuter.

\mr{नियंत्रणात} \mr{येणे}, to come under control.

\begin{cousinweb}
\textbf{\mr{नियंत्रण}} (\emph{niyantraṇ}) is a tatsama, taken from Sanskrit \emph{niyantraṇa}, \textquotedblleft{}restraining, governing\textquotedblright{} (Monier-Williams).
\end{cousinweb}

\begin{grammarlens}[title={What you've built}]
One more word for saying what happens if.
\end{grammarlens}

\subsection*{Your turn}
\begin{itemize}
\raggedright
  \item \emph{Say it:} \emph{niyantraṇ}
  \item \emph{Say it:} \emph{niyantraṇ}, once more
  \item \emph{Say it:} \emph{āg lavkarac niyantraṇāt ālī.} — The fire was soon brought under control.
  \item \emph{Recall:} say the Marathi for to issue a warning, then the Marathi for a rulebook, then say \emph{niyantraṇ} again
\end{itemize}

\begin{tcolorbox}[breakable,colback=teal!4,colframe=teal!35!black,title={Before you move on}]
What does \mr{नियंत्रण} mean? (\textquotedblleft{}Control\textquotedblright{}.) Say it once more.
\end{tcolorbox}
\section[\texorpdfstring{\mr{ताकीद} (tākīd)}{ताकीद (tākīd)}]{\mr{ताकीद} (tākīd) — a stern warning; a reprimand}

\label{lesson:MR-C542-takid}



\begin{quote}
Before the new one: say the Marathi for a rulebook, then the Marathi for control.
\end{quote}

\subsection*{You'll want to know: \mr{ताकीद}}
\textbf{\mr{ताकीद}} — \emph{tākīd} — \textquotedblleft{}a stern warning; a reprimand\textquotedblright{}.

\begin{quote}
\mr{उशिरा} \mr{आल्यामुळे} \mr{व्यवस्थापकाने} \mr{त्याला} \mr{ताकीद} \mr{दिली}.
\end{quote}

\begin{quote}
\emph{uśirā ālyāmuḷe vyavasthāpakāne tyālā tākīd dilī.}
\end{quote}

\begin{quote}
Because he came late, the manager gave him a stern warning.
\end{quote}

Feminine.

\mr{ताकीद} \mr{देणे}, to warn sternly.

\begin{cousinweb}
\textbf{\mr{ताकीद}} (\emph{tākīd}) is a loan through Persian usually traced to Arabic \emph{taʾkīd}, \textquotedblleft{}emphasis, insistence\textquotedblright{}.
\end{cousinweb}

\begin{grammarlens}[title={What you've built}]
One more word for saying what happens if.
\end{grammarlens}

\subsection*{Your turn}
\begin{itemize}
\raggedright
  \item \emph{Say it:} \emph{tākīd}
  \item \emph{Say it:} \emph{tākīd}, once more
  \item \emph{Say it:} \emph{uśirā ālyāmuḷe vyavasthāpakāne tyālā tākīd dilī.} — Because he came late, the manager gave him a stern warning.
  \item \emph{Recall:} say the Marathi for a rulebook, then the Marathi for control, then say \emph{tākīd} again
\end{itemize}

\begin{tcolorbox}[breakable,colback=teal!4,colframe=teal!35!black,title={Before you move on}]
What does \mr{ताकीद} mean? (\textquotedblleft{}A stern warning; a reprimand\textquotedblright{}.) Say it once more.
\end{tcolorbox}
\section[\texorpdfstring{\mr{जुमानणे} (jumānṇe)}{जुमानणे (jumānṇe)}]{\mr{जुमानणे} (jumānṇe) — to heed; to care about}

\label{lesson:MR-C542-jumanne}



\begin{quote}
Before the new one: say the Marathi for control, then the Marathi for a stern warning.
\end{quote}

\subsection*{You'll want to know: \mr{जुमानणे}}
\textbf{\mr{जुमानणे}} — \emph{jumānṇe} — \textquotedblleft{}to heed; to care about\textquotedblright{}.

\begin{quote}
\mr{तो} \mr{कोणाचाच} \mr{सल्ला} \mr{जुमानत} \mr{नाही}.
\end{quote}

\begin{quote}
\emph{to koṇācāc sallā jumānat nāhī.}
\end{quote}

\begin{quote}
He heeds nobody's advice.
\end{quote}

Past, with -\mr{ने} (he …): \mr{त्याने} \mr{आदेश} \mr{जुमानला}, \emph{tyāne ādeś jumānlā}.

Mostly with \mr{नाही}: \mr{तो} \mr{कोणालाच} \mr{जुमानत} \mr{नाही}, he heeds no one.

\begin{cousinweb}
\textbf{\mr{जुमानणे}} (\emph{jumānṇe}) is a Marathi verb heard mostly with a negative; its origin is not settled.
\end{cousinweb}

\begin{grammarlens}[title={What you've built}]
One more word for saying what happens if.
\end{grammarlens}

\subsection*{Your turn}
\begin{itemize}
\raggedright
  \item \emph{Say it:} \emph{jumānṇe}
  \item \emph{Say it:} \emph{jumānṇe}, once more
  \item \emph{Say it:} \emph{to koṇācāc sallā jumānat nāhī.} — He heeds nobody's advice.
  \item \emph{Recall:} say the Marathi for control, then the Marathi for a stern warning, then say \emph{jumānṇe} again
\end{itemize}

\begin{tcolorbox}[breakable,colback=teal!4,colframe=teal!35!black,title={Before you move on}]
What does \mr{जुमानणे} mean? (\textquotedblleft{}To heed; to care about\textquotedblright{}.) Say it once more.
\end{tcolorbox}
\section[\texorpdfstring{\mr{झुगारणे} (jhugārṇe)}{झुगारणे (jhugārṇe)}]{\mr{झुगारणे} (jhugārṇe) — to defy; to throw off}

\label{lesson:MR-C542-jhugarne}



\begin{quote}
Before the new one: say the Marathi for a stern warning, then the Marathi for to heed.
\end{quote}

\subsection*{You'll want to know: \mr{झुगारणे}}
\textbf{\mr{झुगारणे}} — \emph{jhugārṇe} — \textquotedblleft{}to defy; to throw off\textquotedblright{}.

\begin{quote}
\mr{काही} \mr{लोकांनी} \mr{नियम} \mr{झुगारले} \mr{आणि} \mr{दंड} \mr{भरला}.
\end{quote}

\begin{quote}
\emph{kāhī lokānnī niyam jhugārle āṇi daṇḍ bharlā.}
\end{quote}

\begin{quote}
Some people defied the rules and paid a fine.
\end{quote}

Past, with -\mr{ने} (he …): \mr{त्याने} \mr{नियम} \mr{झुगारले}, \emph{tyāne niyam jhugārle}.

\begin{cousinweb}
\textbf{\mr{झुगारणे}} (\emph{jhugārṇe}) is a Marathi verb for throwing off a rule or a burden; its origin is not settled.
\end{cousinweb}

\begin{grammarlens}[title={What you've built}]
One more word for saying what happens if.
\end{grammarlens}

\subsection*{Your turn}
\begin{itemize}
\raggedright
  \item \emph{Say it:} \emph{jhugārṇe}
  \item \emph{Say it:} \emph{jhugārṇe}, once more
  \item \emph{Say it:} \emph{kāhī lokānnī niyam jhugārle āṇi daṇḍ bharlā.} — Some people defied the rules and paid a fine.
  \item \emph{Recall:} say the Marathi for a stern warning, then the Marathi for to heed, then say \emph{jhugārṇe} again
\end{itemize}

\begin{tcolorbox}[breakable,colback=teal!4,colframe=teal!35!black,title={Before you move on}]
What does \mr{झुगारणे} mean? (\textquotedblleft{}To defy; to throw off\textquotedblright{}.) Say it once more.
\end{tcolorbox}
